% Options for packages loaded elsewhere
% Options for packages loaded elsewhere
\PassOptionsToPackage{unicode}{hyperref}
\PassOptionsToPackage{hyphens}{url}
\PassOptionsToPackage{dvipsnames,svgnames,x11names}{xcolor}
%
\documentclass[
  11pt,
]{article}
\usepackage{xcolor}
\usepackage[top=0.4in,bottom=0.45in,left=0.65in,right=0.65in,includehead,headheight=52pt,headsep=7pt,includefoot,footskip=24pt]{geometry}
\usepackage{amsmath,amssymb}
\setcounter{secnumdepth}{-\maxdimen} % remove section numbering
\usepackage{iftex}
\ifPDFTeX
  \usepackage[T1]{fontenc}
  \usepackage[utf8]{inputenc}
  \usepackage{textcomp} % provide euro and other symbols
\else % if luatex or xetex
  \usepackage{unicode-math} % this also loads fontspec
  \defaultfontfeatures{Scale=MatchLowercase}
  \defaultfontfeatures[\rmfamily]{Ligatures=TeX,Scale=1}
\fi
\usepackage{lmodern}
\ifPDFTeX\else
  % xetex/luatex font selection
\fi
% Use upquote if available, for straight quotes in verbatim environments
\IfFileExists{upquote.sty}{\usepackage{upquote}}{}
\IfFileExists{microtype.sty}{% use microtype if available
  \usepackage[]{microtype}
  \UseMicrotypeSet[protrusion]{basicmath} % disable protrusion for tt fonts
}{}
\makeatletter
\@ifundefined{KOMAClassName}{% if non-KOMA class
  \IfFileExists{parskip.sty}{%
    \usepackage{parskip}
  }{% else
    \setlength{\parindent}{0pt}
    \setlength{\parskip}{6pt plus 2pt minus 1pt}}
}{% if KOMA class
  \KOMAoptions{parskip=half}}
\makeatother
% Make \paragraph and \subparagraph free-standing
\makeatletter
\ifx\paragraph\undefined\else
  \let\oldparagraph\paragraph
  \renewcommand{\paragraph}{
    \@ifstar
      \xxxParagraphStar
      \xxxParagraphNoStar
  }
  \newcommand{\xxxParagraphStar}[1]{\oldparagraph*{#1}\mbox{}}
  \newcommand{\xxxParagraphNoStar}[1]{\oldparagraph{#1}\mbox{}}
\fi
\ifx\subparagraph\undefined\else
  \let\oldsubparagraph\subparagraph
  \renewcommand{\subparagraph}{
    \@ifstar
      \xxxSubParagraphStar
      \xxxSubParagraphNoStar
  }
  \newcommand{\xxxSubParagraphStar}[1]{\oldsubparagraph*{#1}\mbox{}}
  \newcommand{\xxxSubParagraphNoStar}[1]{\oldsubparagraph{#1}\mbox{}}
\fi
\makeatother


\usepackage{longtable,booktabs,array}
\newcounter{none} % for unnumbered tables
\usepackage{calc} % for calculating minipage widths
% Correct order of tables after \paragraph or \subparagraph
\usepackage{etoolbox}
\makeatletter
\patchcmd\longtable{\par}{\if@noskipsec\mbox{}\fi\par}{}{}
\makeatother
% Allow footnotes in longtable head/foot
\IfFileExists{footnotehyper.sty}{\usepackage{footnotehyper}}{\usepackage{footnote}}
\makesavenoteenv{longtable}
\usepackage{graphicx}
\makeatletter
\newsavebox\pandoc@box
\newcommand*\pandocbounded[1]{% scales image to fit in text height/width
  \sbox\pandoc@box{#1}%
  \Gscale@div\@tempa{\textheight}{\dimexpr\ht\pandoc@box+\dp\pandoc@box\relax}%
  \Gscale@div\@tempb{\linewidth}{\wd\pandoc@box}%
  \ifdim\@tempb\p@<\@tempa\p@\let\@tempa\@tempb\fi% select the smaller of both
  \ifdim\@tempa\p@<\p@\scalebox{\@tempa}{\usebox\pandoc@box}%
  \else\usebox{\pandoc@box}%
  \fi%
}
% Set default figure placement to htbp
\def\fps@figure{htbp}
\makeatother





\setlength{\emergencystretch}{3em} % prevent overfull lines

\providecommand{\tightlist}{%
  \setlength{\itemsep}{0pt}\setlength{\parskip}{0pt}}



 


% =====================================================================
%  _style.tex  ·  Design system for EIL research highlights (v4 layout)
%  Edit here to restyle ALL research highlights. Included via the .qmd YAML.
%  Matches the canonical single-column v4 template; the masthead and footer
%  are set per-document via fancyhdr in each .qmd's include-in-header.
% =====================================================================

% ---- Colours (exact hex) --------------------------------------------
\definecolor{ink}{HTML}{1D2620}        % headings, title, section heads, rule
\definecolor{body}{HTML}{33352C}       % main body paragraphs
\definecolor{accentred}{HTML}{641111}  % RESEARCH HIGHLIGHT label, links, fig title
\definecolor{boxred}{HTML}{8C2826}     % Bottom Line box fill
\definecolor{boxtext}{HTML}{EEF1EA}    % Bottom Line body text (off-white)
\definecolor{muted}{HTML}{50534A}      % figure source line
\definecolor{cite}{HTML}{6E6D61}       % citation line
\definecolor{faint}{HTML}{9A978A}      % footer
\definecolor{warmrule}{HTML}{D8D2C2}   % thin warm-grey dividers
\definecolor{paper}{HTML}{FAF8F1}      % warm page background

% \pagecolor{paper}                    % uncomment for warm off-white background

% ---- Fonts (Source Sans 3 / Source Serif 4) -------------------------
\usepackage[T1]{fontenc}
\usepackage[default]{sourcesanspro}
\usepackage{sourceserifpro}
\renewcommand{\familydefault}{\sfdefault}

% ---- Page -----------------------------------------------------------
\pagestyle{empty}
\setlength{\parindent}{0pt}
\setlength{\parskip}{2pt}
% Default body size/colour. The v4 .qmd overrides this to 10.5/13.5pt.
\AtBeginDocument{\color{body}\fontsize{9.4pt}{13.6pt}\selectfont}

% ---- Packages -------------------------------------------------------
\usepackage{microtype}
\usepackage{soul}
\setul{2pt}{.5pt}
\usepackage{tcolorbox}
\tcbuselibrary{skins}
\usepackage[colorlinks=true, urlcolor=accentred, linkcolor=accentred, pdfborder={0 0 0}]{hyperref}

% =====================================================================
%  MACROS
% =====================================================================

% --- Title + citation (citation has a thin warm divider beneath) ------
\newcommand{\brieftitle}[1]{%
  {\color{ink}\fontsize{18pt}{19.5pt}\selectfont\bfseries\textls[-12]{#1}\par}\vspace{3pt}}
\newcommand{\briefcitation}[1]{%
  {\color{cite}\fontsize{9.4pt}{12pt}\selectfont #1\par}\vspace{3pt}%
  \noindent\textcolor{warmrule}{\rule{\linewidth}{0.7pt}}}

% --- Section heads (uppercase, ink) ----------------------------------
\newcommand{\sectionhead}[1]{%
  \vspace{2pt}{\color{ink}\bfseries\fontsize{9.75pt}{12pt}\selectfont
  \textls[20]{\MakeUppercase{#1}}\par}\vspace{1pt}}

% --- Sub-heads (small uppercase) — Takeaway subheads -----------------
% The v4 .qmd overrides this to the brand maroon via \renewcommand.
\newcommand{\subhead}[1]{%
  {\color{ink}\bfseries\fontsize{8.25pt}{10pt}\selectfont
  \textls[50]{\MakeUppercase{#1}}\par}\vspace{2pt}}

% --- Figure title + source line --------------------------------------
\newcommand{\figtitle}[2]{%
  {\bfseries\color{accentred}\fontsize{7.1pt}{9pt}\selectfont\textls[40]{\MakeUppercase{#1}}}\par%
  {\color{muted}\fontsize{7.1pt}{9pt}\selectfont #2}\par\vspace{4pt}}

% =====================================================================
%  COLOURED BOXES
% =====================================================================

% --- Bottom Line box (border-radius 6px -> arc=4pt) ------------------
\newtcolorbox{bottomline}{%
  enhanced, arc=4pt, outer arc=4pt, boxrule=0pt,
  colback=boxred, coltext=boxtext,
  left=7pt, right=7pt, top=7pt, bottom=7pt,
  before upper={\fontsize{9.4pt}{13.6pt}\selectfont}}
% --- Running masthead + footer (logo, label, rule, page number) ---
\usepackage{fancyhdr}
\renewcommand{\headrulewidth}{0pt}
\renewcommand{\footrulewidth}{0pt}
\fancypagestyle{brief}{%
  \fancyhf{}%
  \fancyhead[L]{%
    \begin{minipage}[c]{0.62\textwidth}\raggedright\color{ink}
\includegraphics[height=0.45in]{../../../formats/logos/eil-logo-maroon.png}\end{minipage}\hfill%
    \begin{minipage}[c]{0.36\textwidth}\raggedleft\fontsize{8.25pt}{10pt}\selectfont\bfseries\color{accentred}\textls[200]{\MakeUppercase{Research Highlight}}\end{minipage}%
    \par\vspace{2pt}%
    \noindent\textcolor{ink}{\rule{\textwidth}{2.2pt}}%
  }%
  \fancyfoot[C]{%
    \color{warmrule}\rule{\textwidth}{0.7pt}\par\vspace{3pt}%
    \hbox to\textwidth{%
      \color{faint}\fontsize{7.4pt}{9pt}\selectfont
      For full paper:\hspace{0.4em}\href{https://researchonline.lse.ac.uk/id/eprint/126795}{\color{faint}\underline{\textls[40]{researchonline.lse.ac.uk/id/eprint/126795}}}%
      \hfill
      Page \thepage%
    }%
  }%
}
\pagestyle{brief}
% Takeaways subheads in the brand maroon (logo / Research Highlight)
\renewcommand{\subhead}[1]{%
  {\color{accentred}\bfseries\fontsize{8.25pt}{10pt}\selectfont
  \textls[50]{\MakeUppercase{#1}}\par}\vspace{2pt}}
\makeatletter
\@ifpackageloaded{caption}{}{\usepackage{caption}}
\AtBeginDocument{%
\ifdefined\contentsname
  \renewcommand*\contentsname{Table of contents}
\else
  \newcommand\contentsname{Table of contents}
\fi
\ifdefined\listfigurename
  \renewcommand*\listfigurename{List of Figures}
\else
  \newcommand\listfigurename{List of Figures}
\fi
\ifdefined\listtablename
  \renewcommand*\listtablename{List of Tables}
\else
  \newcommand\listtablename{List of Tables}
\fi
\ifdefined\figurename
  \renewcommand*\figurename{Figure}
\else
  \newcommand\figurename{Figure}
\fi
\ifdefined\tablename
  \renewcommand*\tablename{Table}
\else
  \newcommand\tablename{Table}
\fi
}
\@ifpackageloaded{float}{}{\usepackage{float}}
\floatstyle{ruled}
\@ifundefined{c@chapter}{\newfloat{codelisting}{h}{lop}}{\newfloat{codelisting}{h}{lop}[chapter]}
\floatname{codelisting}{Listing}
\newcommand*\listoflistings{\listof{codelisting}{List of Listings}}
\makeatother
\makeatletter
\makeatother
\makeatletter
\@ifpackageloaded{caption}{}{\usepackage{caption}}
\@ifpackageloaded{subcaption}{}{\usepackage{subcaption}}
\makeatother
\usepackage{bookmark}
\IfFileExists{xurl.sty}{\usepackage{xurl}}{} % add URL line breaks if available
\urlstyle{same}
\hypersetup{
  colorlinks=true,
  linkcolor={blue},
  filecolor={Maroon},
  citecolor={Blue},
  urlcolor={Blue},
  pdfcreator={LaTeX via pandoc}}


\author{}
\date{}
\begin{document}


% ---------------------------------------------------------------------
% TITLE + CITATION  (masthead repeats in the page header)
% ---------------------------------------------------------------------
\brieftitle{Coal's Decline Imposed Large, Lasting Costs on Its Workers}
\briefcitation{Lessons from ``Transitional Costs and the Decline of Coal: Worker-Level Evidence'' by Jonathan Colmer, Eleanor Krause, Eva Lyubich \& John Voorheis, Working Paper, 2025.}

\vspace{-2pt}

% override body font size for this document
\fontsize{10.5pt}{13.5pt}\selectfont

% ---------------------------------------------------------------------
% BOTTOM LINE  (~40–50 words)
% ---------------------------------------------------------------------
\begin{bottomline}
{\fontsize{11pt}{15pt}\selectfont
\textbf{When the U.S. coal industry shrank by half, the workers it left behind never caught back up.} The most coal-dependent workers lost more than a full year's earnings over the following decade --- and neither switching industries nor moving to a new region did much to soften the blow.
}
\end{bottomline}

\vspace{6pt}

% ---------------------------------------------------------------------
% BACKGROUND  (~120 words, 2 paragraphs)
% ---------------------------------------------------------------------
\sectionhead{Background}
\noindent U.S. coal mining employment has fallen more than half since 2011, as cheap natural gas from the fracking boom and falling renewable energy costs pushed coal out of the power sector. Earlier research shows that coal \textit{communities} lost jobs, wages, and tax revenue. However, aggregate numbers hide how individual workers fared. Some may have moved on quickly to new jobs, while others may have borne heavy \textbf{transitional costs} --- the lasting losses to earnings and employment that workers absorb when an industry contracts. Which of these patterns held for coal workers had never been measured at the worker level.

\noindent The research team assembled administrative records covering nearly every person who worked in U.S. coal mining before, during, and after coal's 2011 peak. This let them follow each worker's earnings, employment, and location year by year, even through years spent outside coal or out of work entirely. They could then compare those paths with otherwise similar workers who depended less on coal for their livelihoods.

\vspace{6pt}

% ---------------------------------------------------------------------
% WHAT THE STUDY FOUND  (~150 words + figure)
% ---------------------------------------------------------------------
\sectionhead{What the study found}
% ---------------------------------------------------------------------
% FIGURE — recreation of the paper's Fig. 7(a); built by code/make-fig1.R
% ---------------------------------------------------------------------
\begin{center}
\begin{minipage}{0.82\linewidth}
\figtitle{Coal workers' earnings fell further behind comparable workers every year}{Source: recreated from Colmer et al.\ (2025), Fig.\ 7a.}
\includegraphics[width=\linewidth]{../data-viz/figures/fig1.png}
\par\vspace{3pt}
{\color{faint}\fontsize{7.1pt}{9.5pt}\selectfont \textit{Note:} Each point is the cumulative earnings gap between workers whose livelihoods depended most on coal and otherwise similar workers who depended on it less, measured in years of pre-decline pay. The two groups track together until coal's 2011 peak, then steadily diverge --- by 2019 the gap reaches 1.6 years' worth of earnings. Values digitized from the published figure.\par}
\end{minipage}
\end{center}

\vspace{6pt}

\newpage

\noindent Measuring what coal's decline cost workers means separating its effects from everything else shaping a career. The researchers compared the most coal-dependent workers with other workers, evening out gaps in age, education, earnings, and location, so that whatever difference emerged afterward could be chalked up to coal's decline. The collapse in coal demand was driven by national energy prices rather than any one worker's choices, so the comparison offers a clean picture of coal's effect. If coal's decline had left workers unharmed --- or had weighed on both groups equally --- no earnings gap should have opened between them.

\noindent This is not what the researchers found. Between 2012 and 2019, \textbf{the most coal-dependent workers lost cumulative earnings equal to roughly 1.6 times a full pre-shock year of pay}. The losses came through two channels at once: compared with similar workers, they spent more time out of work \textit{and} earned less in the years they did work. \textbf{A decade on, coal's decline had carved a deep and persistent hole in workers' earnings.}

\vspace{8pt}

% ---------------------------------------------------------------------
% TAKEAWAYS  (3 points, ~50 words each)
% ---------------------------------------------------------------------
\sectionhead{Takeaways}
\subhead{Nowhere to go}
\noindent Workers couldn't simply move or retrain their way out. Coal employment is concentrated in a few isolated regions where entire local economies lean on the industry, so nearby alternatives were scarce. Those who left coal earned about 30 percent less, and relocating to a new labor market barely helped. Where workers displaced by factory closings often find their way into new jobs, a striking share of coal workers stopped working altogether --- \textbf{a sign their skills didn't transfer to the jobs around them}.

\vspace{5pt}
\subhead{Leaning on the safety net}
\noindent As earnings fell, coal-dependent workers leaned on government transfers --- especially Social Security Disability Insurance (SSDI) --- and drew down retirement savings earlier. These payments cushioned the blow but \textbf{replaced only a fraction of what workers lost}, a second-best safety net standing in for the jobs and wages that disappeared.

\vspace{5pt}
\subhead{A preview of the energy transition}
\noindent Coal's decline is an early look at what shrinking demand for fossil fuels could mean for workers in concentrated, high-wage industries. Coal is just one industry, and its geography is unusually isolating, so its exact path won't repeat everywhere. Even so, the warning is clear --- \textbf{where specialized skills don't transfer, the costs of an energy transition land hard on workers, and existing safety nets catch only part of the fall}.




\end{document}
